\chapter*{Abstract}\label{cha:abstract}
\addcontentsline{toc}{chapter}{Abstract}


KI eröffnet der Luftfahrt neue Möglichkeiten, trifft dort 
jedoch auf besonders hohe Erwartungen an Zuverlässigkeit, Robustheit und nachvollziehbare 
Absicherung. Am Beispiel von Luftfahrtradarsystemen wird der aktuelle Diskussionsstand zur 
Sicherheit von KI eingeordnet: Im Zentrum steht, welche Absicherungsansätze gefordert 
werden und in welchem Maß diese durch heutige Praktiken abgedeckt erscheinen. Als 
Bezugsrahmen dienen hierbei etablierte Sicherheitsstandards – u. a. die ISO 8800 – 
sowie konzeptionelle Leitlinien zuständiger Luftfahrtbehörden.

Die Arbeit stützt sich auf eine systematische Literaturrecherche, definiert die 
Grundlagen zu KI, Radar und Sicherheit und erläutert ausgewählte KI-Modelle mit Blick auf 
Erklärbarkeit und typische Fehlerbilder. Ergänzend werden Praxisbeobachtungen und Berichte 
zu Problemen beim KI-Einsatz zusammengetragen und aufbereitet. Darauf aufbauend werden typische Anforderungen 
und Szenarien für Radarsysteme herausgearbeitet und bestehende AI-Safety-Ansätze entlang 
der ISO 8800 analysiert, deren Übertragbarkeit auf den Radar-Kontext 
bewertet sowie potenzielle Sicherheitslücken und Risiken identifiziert. Aus den Ergebnissen 
werden anwendungsbezogene Empfehlungen abgeleitet. 
Auf Basis eines ausgewählten Beispielmodells wurden mögliche Risiken identifiziert. Diese wurden verglichen 
mit einer ergänzenden Analyse zu heutigen Standards.

Ergebnis: Die identifizierten Risiken sind bereits erfasst und methodisch abgedeckt.

\newpage
